\chapter{General conclusion}
The primary aim of this study was to deepen the understanding of subject focus and its realization in Spanish and French. To achieve this, I adopted a quantitative variationist approach, treating subject focus as a linguistic variable and its various realizations as optional variants. Although Spanish and French belong to the same linguistic family, they are said to display significant differences in how they encode focus, particularly subject focus. Theoretical accounts have traditionally claimed that French expresses subject focus primarily through constructional strategies such as \textit{c'est} clefts, while Spanish does so through postverbal subjects. These accounts often suggest a one-to-one correlation between word order patterns and discourse-functional motivations, thereby downplaying the potential for variation (for Spanish: \cite{zubizarreta1998prosody, costa2001postverbal, samek2001crosslinguistic, gutierrez2013structural}, \textit{inter alia}. For French: \cite{prince1978comparison, lambrecht1996information, belletti2005answering}, \textit{inter alia}).

In contrast, empirical studies have indicated that subject focus in both languages can be realized through various competing strategies, pointing to the presence of variation in this phenomenon (for Spanish: \cite{gabriel2007fokus, gabriel2010focus, heidinger2013information, uth2014spanish, cabrera1999funciones, feldhausen2015oraciones}, \textit{inter alia}. For French: \cite{dufter2009focus, fery2001focus, choi2005gens, lahousse2006, Lahousse2011, karssenberg2017oiseaux, Karssenberg2018}, \textit{inter alia}). However, both theoretical and empirical studies have largely concentrated on full sentences, neglecting cases where subject focus is encoded by elliptical structures, such as reduced clefts and fragments.

This book seeks to address this gap by conducting a corpus-based study of spontaneous dialogues in both languages. By identifying both full and elliptical strategies that speakers employ to express subject focus, this study provides new insights into the factors governing their alternation. Below, I offer a summary of the key contributions of this book.

\section{Frequency of subject focus-marking strategies} %(Chapter 3)


In Chapter 3, I delineated the envelope of variation by selecting criteria that defined the variable of subject focus. Using the Question Under Discussion (QUD) approach to focus \parencite{roberts1996information, buring2003d, beaver2009sense}, I extracted various focus-marking strategies for both Spanish and French and annotated them based on the type of focus they conveyed, distinguishing between contrastive and non-contrastive readings. In both languages, multiple strategies were available to express the same function, though they varied in frequency.

In the Spanish corpus, speakers used postverbal subjects, fragments, \textit{in situ} focus, clefts, and reduced clefts. By contrast, in the French dataset, speakers employed clefts, reduced clefts, \textit{in situ} focus, fragments, and dislocations. Notably, no cases of postverbal subjects appeared in the French corpus, and no instances of dislocations were found in the Spanish corpus.

Regarding focus type, the majority of cases in both languages were classified as information focus, i.e., focus expressing non-presupposed information. Contrastive foci, however, were rare in both datasets, aligning with previous accounts that characterize contrast as a more marked pragmatic function and, therefore, less frequent than simply conveying new information \parencite{zimmermann2011focus}.

Regarding the frequencies of the various strategies, postverbal subjects (n=167) and cleft sentences (n=220) emerged as the most common response types in Spanish and French, respectively. Notably, the second most frequent strategy in both languages was an elliptical one: Spanish speakers produced 102 fragment answers, while French speakers produced 137 reduced clefts. This suggests that even when employing elliptical answers, the two languages exhibit significant differences.
To explain this observation, I posited the hypothesis that speakers in both languages tend to `copy” the focused part of the most frequently used full strategy and elide the background proposition for discourse-related reasons.

Furthermore, the results from the Spanish corpus indicated that in situ structures and clefts are not exclusively associated with a contrastive interpretation; they can also express information focus. This finding challenges previous accounts that asserted a necessary connection between the use of cleft sentences or \textit{in situ} strategies in Spanish and a contrastive or corrective interpretation \parencite{lambrecht1996information, zubizarreta1998prosody}. Instead, it aligns with the analyses by \textcite{fery2001focus} and \textcite{gabriel2007fokus}, who explored prosodic focus in Spanish, as well as with \textcite{feldhausen2015oraciones}, who examined peninsular Spanish clefts.

Overall, this chapter has confirmed the assertion that focused subjects in French and Spanish can be expressed through various, alternating strategies, highlighting the absence of a one-to-one correlation between form and function. Instead, multiple syntactic structures can accommodate different readings, with the appropriate interpretation emerging solely from the discourse contexts in which these structures are situated. Additionally, the corpus study conducted in this chapter identified and analysed different elliptical structures that compete with full forms in expressing the same information-structural function. The investigation of the contexts in which these structures occur and are preferred over their full counterparts was the focus of Chapter 4.




\section{The alternation between full and elliptical answers} %(Chapter 4)


In Chapter 4, I investigated the factors underlying the alternation between full and elliptical strategies that emerged from the corpus analysis in Chapter 3. Previous literature had suggested that reduced clefts in French are particularly well-suited as responses to immediately preceding questions \parencite{belletti2008answering, katz2000categories}. However, this assertion had never been empirically tested until now.

In my analysis, I postulated the hypothesis that while elliptical answers (such as reduced clefts and fragments) are more suitable for immediate QUDs, full strategies are predominantly used in response to distant or implicit QUDs — what I called the ``immediacy hypothesis”. This hypothesis is grounded in the observation that when the QUD is immediate, the background information from the question is highly activated in the minds of both the speaker and the listener. In such cases, repeating the entire background in the answer can feel unnatural. Instead, opting for an elliptical response, where only the focused part remains, is more appropriate. Conversely, when the QUD is distant or is implicit, the activation of background material decreases, prompting a cooperative speaker to favour a full structure.
To test this hypothesis empirically, I annotated each QUD preceding the answer that expressed subject focus and distinguished between three levels of annotation: i) Immediate QUD, ii) Distant QUD, iii) Implicit QUD. The results from both languages supported this hypothesis, although certain tendencies were more pronounced in French than in Spanish. In both languages, full structures were more frequent when preceded by distant or implicit QUDs. However, in the immediate condition, French speakers showed a strong inclination to respond with elliptical strategies, whereas this tendency was considerably weaker in the Spanish dataset.

In both languages, syntactic structures such as clefts or i\textit{n situ} responses, which contain both the focused element and the background proposition, tend to be more frequent when the triggering QUD is distant or implicit. In the French dataset, however, elliptical structures — like reduced clefts or fragments — where only the focused element remains while the background proposition is omitted, are more commonly found as answers to immediate QUDs.

To explain this discrepancy between the two languages in the immediate condition, I proposed a reasoning based on syntactic complexity. As established in Chapter 3, the most frequent full answer type in French is a cleft sentence, while in Spanish, the corresponding full strategy appears to be a postverbal subject. These two structures differ inherently at the syntactic level: postverbal subjects result from a reordering of constituents in the sentence, whereas clefts are bi-clausal structures comprised of a main clause and a subordinate clause. This bi-clausal configuration renders cleft sentences inherently heavier and more complex than postverbal subjects.
Moreover, the analysis of verb types in Chapter 5 indicated that most postverbal subjects involve unaccusative verbs, which means that these sentences consist of just a subject and a verb, lacking any additional arguments.

Based on these observations, I proposed that repeating the entire background proposition through a separate relative clause in a cleft is perceived as heavier and less felicitous by speakers than merely repeating a single verb in a postverbal subject. This perception may account for why Spanish speakers did not exhibit a significant preference for either a full or an elliptical answer in response to immediately preceding questions.

This chapter has made several innovative contributions. First, it has operationalized the level of activation of the background proposition in dialogical contexts by distinguishing various levels of immediacy in the preceding QUD. Second, it has explained the occurrence of elliptical answer types and their alternation with full sentences when expressing the same IS function. Third, it has provided evidence for the Principle of Communicative Efficiency \parencite{Levshina_2022}, which posits that when sufficient contextual cues are available for successful understanding, speakers prefer to convey less. Conversely, they tend to express more when such cues are less clear.



\section{Pragmatic and semantic factors in syntactic variation} %(Chapter 5)


In Chapter 5, I examined the alternation between postverbal and preverbal subjects in Spanish, as well as the various cleft formats in French. Previous literature has suggested that the choice between SV and VS order in contexts where the subject is in focus may be influenced by factors such as the givenness of the subject constituent and the type of verb used \parencite{gabriel2007fokus}. In French, the optionality between \textit{c'est} and \textit{il y a} clefts — when expressing the same information structure (IS) function — has often been linked to the type of focus. Specifically, it has been argued that while \textit{c'est} clefts typically convey an exhaustive reading of the focused constituent, \textit{il y a} clefts do so only marginally \parencite{de2015status, karssenberg2017oiseaux}.

In this chapter's analysis, I applied the same predictors to both languages. Specifically, I examined whether the verb’s argument structure, the degree of givenness of the subject constituent, and exhaustivity could explain the syntactic variation observed in the two languages. First, I conducted descriptive and inferential statistics, exploring the main effect of each factor on the different syntactic structures, followed by an investigation of the interaction between the three predictors. To determine the significance of these results, I performed inferential statistics using mixed-effects modelling, with the interview number included as a random effect.

The results revealed that factors such as verb type and givenness align with specific syntactic structures in both languages. Notably, the combination of unaccusative verbs and new subjects correlates strongly with Spanish postverbal subjects and French \textit{il y a} clefts. In contrast, transitive verbs and given subjects are more frequently associated with French \textit{c'est} clefts and, to a lesser extent, with Spanish \textit{in situ} foci. Moreover, the interaction of these two factors appears to reinforce this correlation.

This finding supports the claims made by \textcite{du2003preferred} and the Preferred Argument Structure (PAS), which had previously highlighted the link between transitive verbs and given subjects, as well as unaccusative verbs with new subjects. Additionally, prior research on the VS order had already indicated that this word order is more commonly found with unaccusatives (e.g., \cite{burzio1986italian} for Italian; \cite{gabriel2007fokus} for empirical evidence on Spanish VS). However, the role of argument structure in relation to French clefts had not been systematically explored until now. On the one hand, the corpus study presented in this chapter confirms claims previously made in the literature. On the other hand, it provides new insights into the factors that may drive the optionality of different cleft formats when they compete for the same discursive function.

Regarding the role of exhaustivity, significance tests revealed that this factor does not yield significant results in Spanish. However, in French, there is a noticeable association between \textit{c'est} clefts and exhaustive focus, and \textit{il y a} clefts with a non-exhaustive focus type. Nevertheless, the data shows that this correlation is not deterministic, as both cleft formats can be found with either reading. This suggests that exhaustivity is not semantically encoded in the syntax of these constructions but is pragmatically conveyed, aligning with the views of \textcite{horn1981exhaustiveness, dufter2009clefting, karssenberg2017oiseaux, destruel2018interpretation}, \textit{inter alia}. Additionally, some \textit{c'est} clefts occur with exclusive particles like \textit{seulement} (`only'), while some \textit{il y a} clefts co-occur with additive particles such as \textit{aussi} (`too'). If exhaustivity (or its absence) were semantically encoded in clefts, these particles would be incompatible with these constructions and thus not present.

While this book has provided valuable insights and addressed both longstanding and emerging questions, several unresolved issues remain. For instance, certain syntactic structures, such as French pseudo-clefts and dislocations, were too infrequent to be included in the statistical analysis, making generalizations impossible. Additionally, the interaction between the examined factors often resulted in a reduced number of occurrences, limiting the interpretability of the findings. A forced-choice task or an acceptability judgment test, conducted in a more controlled experimental setting where these factors are deliberately combined or mismatched, would offer more definitive evidence of the behaviour of the syntactic constructions under investigation.

Furthermore, the role of prosody was not considered in this study. The \textit{in situ} strategies in both languages, where focus is conveyed by assigning intonational prominence to the subject, clearly call for a prosodic analysis. Such an analysis would offer deeper insights into how the two languages differ in their use of prosody for discourse-related functions.
In short, much work remains to be done before we can fully understand how pragmatic functions are mapped onto syntactic structures. It is my hope that this book will inspire future research on these complex and fascinating phenomena.
